\subsection{Overlapping minimal siphons support distinct invasion thresholds and interior cusps}

\paragraph{Reading siphons from the production.}
\beR[Gaussian and logistic responses have no siphons]
[no siphons][KBnosiph]
Gaussian and logistic CDFs are strictly positive at every finite argument.
They therefore exclude nonempty siphons whenever $a_i(0)>0$ for all $i$.
Zero-onset responses can admit siphons, depending on $h$, $P$ and $V$;
irreducibility of the interaction graph alone does not settle the question.
\eeR
\beR[Graph rule for siphons]
[graph rule][KBsiphgraph]
For zero stimulus, linear nonnegative $P_i$ and such zero-onset
responses, Proposition~\ref{p:InvF} gives a graph description;
with nondiagonal $V$, include transfer edges $j\to i$ whenever
$-v_{ij}>0$, as well as the positive edges of $W$. Siphons are the
sets containing every predecessor of each of their indices, and
minimal siphons are the source strongly connected components of
this combined graph. This graph rule requires a response positive
at every positive argument; a response with an additional flat zero
interval on a finite box requires \eqref{eq:KB-response-zero-test}.
\eeR
\beR[The response has a definite order along transverse rays]
[transverse order][KBtransOrder]
For the threshold model on a siphon face, fix $h_i=0$, a positive resident state $y$,
and a nonzero polynomial $P_i$ with $P_i(0,y)=0$. Put
\[
 d_i=\min_{a_{i\nu}>0}\sum_{j\in \Sigma}\nu_j\ge1,\qquad
 \Phi_i(t)=c_it^{k_i}+o(t^{k_i})\quad(t\downarrow0),\quad c_i,k_i>0.
\]
For each fixed $\zeta\gg0$, the leading coefficient
$p_i(\zeta,y)>0$ satisfies
\[
 P_i(\varepsilon\zeta,y)=\varepsilon^{d_i}p_i(\zeta,y)
 +o(\varepsilon^{d_i}),\qquad
 \Pi_i(\varepsilon\zeta,y)
 =b_ic_ip_i(\zeta,y)^{k_i}\varepsilon^{d_ik_i}
 +o(\varepsilon^{d_ik_i}).
\]
Thus the response has order $d_ik_i$ along positive transverse rays.
For the model with factors $a_i$ this order also holds if $a_i(0)>0$;
a vanishing $a_i(0)$ can increase it.
The density has order $k_i-1$ only if additionally
$\Phi_i'(t)=k_ic_it^{k_i-1}+o(t^{k_i-1})$; this does not follow
by differentiating the response expansion.
Grouping monomials by their degree in $z$ gives the leading
term; its coefficient is positive for $\zeta,y\gg0$.
Substitution into the response expansion gives the displayed
production expansion. Multiplication by $a_i(0)+o(1)$ proves the
assertion for the factors $a_i$ when $a_i(0)>0$. No differentiation of an
$o(\cdot)$ remainder is used.
\eeR

\beXa[Minimal siphons can overlap for higher-order $P_i$]
[overlapping minimal siphons][KBoverlap]
On $[0,1]^3$, let $V=I$, $h_i=0$, and
\[
 P_1(x)=x_2,\qquad P_2(x)=x_1x_3,\qquad P_3(x)=x_2,
 \qquad \Pi_i(x)=\Phi_i(P_i(x)),
\]
where $0\le \Phi_i\le1$, $\Phi_i(0)=0$ and $\Phi_i(q)>0$ for $q>0$.
The minimal siphons are $\{1,2\}$ and $\{2,3\}$.
Indeed, absence of $1$ or $3$ forces inclusion of $2$, and absence
of $2$ must be accompanied by absence of $1$ or $3$ to suppress
$x_1x_3$. Neither singleton is a siphon, and $\{1,3\}$ is not one.
Sharing a response on both minimal siphons forces
$\Phi_1=\Phi_2=\Phi_3$. Thus minimal siphons do not in general form a
partition. Smooth admissible responses exist, for example
$\Phi_i(q)=\exp(-1/q)$ for $q>0$ and zero for $q\le0$.
\eeXa

\paragraph{Why Example~\ref{e:KBoverlap} has no nontrivial boundary residents.}
Consider first the unmodified model of Example~\ref{e:KBoverlap},
\[
 \dot x_1=-x_1+\Phi(x_2),\qquad
 \dot x_2=-x_2+\Phi(x_1x_3),\qquad
 \dot x_3=-x_3+\Phi(x_2),
\]
where $\Phi(0)=0$ and $\Phi(q)>0$ for $q>0$. Its minimal siphons
are $\Sigma_{12}=\{1,2\}$ and $\Sigma_{23}=\{2,3\}$.
On $x_1=x_2=0$ the resident equation is $\dot x_3=-x_3$;
on $x_2=x_3=0$ it is $\dot x_1=-x_1$.
In fact, the origin is its only boundary equilibrium: at an
equilibrium, $x_1=0$ or $x_3=0$ forces $x_2=0$ and then all
coordinates to vanish, while $x_2=0$ forces $x_1x_3=0$.
For an odds response with $\Phi'(0+)=a_1$, its Jacobian at the
origin is
\[
 J(0)=\begin{pmatrix}
 -1&a_1&0\\ 0&-1&0\\ 0&a_1&-1
 \end{pmatrix},\qquad \sigma(J(0))=\{-1,-1,-1\}.
\]
The transverse blocks at resident levels $x_3=u$ and $x_1=u$
are, respectively,
\[
 J_{12}^\perp(u)=\begin{pmatrix}-1&a_1\\a_1u&-1\end{pmatrix},
 \qquad
 J_{23}^\perp(u)=\begin{pmatrix}-1&a_1u\\a_1&-1\end{pmatrix}.
\]
Both have spectral abscissa $-1+a_1\sqrt{u}$, but the
resident satisfies $\dot u=-u$, so $u>0$ is not a resident
equilibrium. Thus the overlap is genuine, while the original
model has no nontrivial boundary equilibrium at which to
compare invasion thresholds. In particular, varying the onset
slope does not change its boundary-equilibrium invasion
eigenvalues.

We now add self-feedback without changing the two minimal
siphons. Use the quartic odds response $\Phi_a$ of
equation~\eqref{eq:oddsFamily}, with
\[
 \phi_a(q)=a_1q+a_2q^2+\frac{128}{343}q^3
                    +\frac{1024}{343}q^4,\qquad
 a^0=\left(\frac{342}{343},\frac{368}{343}\right),
\]
and let $r=1/8$, $s=1/2$. By Proposition~\ref{p:Qua},
$-u+\Phi_{a^0}(u)$ has stable roots $0,s$ and a nondegenerate
cusp root $r$. For $a=(a_1,a_2)$ near $a^0$, denote by
$s_a$ the unique root near $s$; it depends smoothly on $a$,
and
\[
 s_a>0,\qquad \lambda_s(a):=\Phi_a'(s_a)-1<0,
 \qquad a_1<1.
\]


\beT[Three-population odds networks with overlapping minimal siphons admit distinct invasion thresholds, multiple stable equilibria, and an interior cusp]\label{t:KBoverlapCusp}
Consider, on $[0,1]^3$, the cooperative threshold network
\begin{equation}\label{eq:KBoverlapCusp}
\begin{aligned}
 \dot x_1&=-x_1+\Phi_a(x_1+\alpha x_2),\\
 \dot x_2&=-x_2+\Phi_a(x_2+\beta x_1x_3),\\
 \dot x_3&=-x_3+\Phi_a(x_3+\gamma x_2),
\end{aligned}
\qquad \alpha,\beta,\gamma>0.
\end{equation}
Keep the two higher coefficients of $\phi_a$ fixed as above.
Then the following conclusions hold.
\begin{enumerate}
\item The minimal siphons are exactly
$\Sigma_{12}=\{1,2\}$ and $\Sigma_{23}=\{2,3\}$;
the only other nonempty siphon is $\{1,2,3\}$.
Every equilibrium has support either $\varnothing$,
$\{1\}$, $\{3\}$, or $\{1,2,3\}$.
For $x\gg0$ the Jacobian is irreducible and Metzler.

\item Besides $E_0=0$, the equilibria
\[
 E_{12}=(0,0,s_a),\qquad E_{23}=(s_a,0,0)
\]
lie on the respective proper siphon faces. The tangential
eigenvalue at each is $\lambda_s(a)<0$. Their transverse
matrices, with absent variables ordered as $(x_1,x_2)$ and
$(x_2,x_3)$, respectively, are
\begin{equation}\label{eq:KBoverlapJtra}
 J_{12}^\perp=\begin{pmatrix}
 a_1-1&a_1\alpha\\ a_1\beta s_a&a_1-1
 \end{pmatrix},\qquad
 J_{23}^\perp=\begin{pmatrix}
 a_1-1&a_1\beta s_a\\ a_1\gamma&a_1-1
 \end{pmatrix}.
\end{equation}
Consequently the two face reproduction numbers and spectral
abscissae are exactly
\begin{equation}\label{eq:KBoverlapR}
\begin{aligned}
 R_{\Sigma_{12}}&=a_1\bigl(1+\sqrt{\alpha\beta s_a}\bigr),&
 s(J_{12}^\perp)&=R_{\Sigma_{12}}-1,\\
 R_{\Sigma_{23}}&=a_1\bigl(1+\sqrt{\gamma\beta s_a}\bigr),&
 s(J_{23}^\perp)&=R_{\Sigma_{23}}-1.
\end{aligned}
\end{equation}
In particular $E_{12}$ and $E_{23}$ are locally
exponentially stable precisely when their respective
reproduction numbers are less than one. With $a$ and
$\beta>0$ fixed, the two thresholds can be varied
independently by $\alpha$ and $\gamma$; they differ whenever
$\alpha\ne\gamma$. For example, the critical values are
\[
 \alpha_c=\gamma_c=
 \frac{(1-a_1)^2}{a_1^2\beta s_a}.
\]
These exact boundary-threshold formulas do not presume the
existence of an interior cusp at the critical couplings.

\item There is $\delta>0$ and a smooth map
\[
 (\alpha,\beta,\gamma)\longmapsto
 \bigl(a_1(\alpha,\beta,\gamma),
       a_2(\alpha,\beta,\gamma),
       x_c(\alpha,\beta,\gamma)\bigr)
\]
defined for $0\le\alpha,\beta,\gamma<\delta$,
with $(a_1,a_2,x_c)=(342/343,368/343,r)$ at zero
coupling, such that \eqref{eq:KBoverlapCusp} has a
nondegenerate interior cusp near $(r,s,s)$. Its zero
eigenvalue is simple, its other two eigenvalues have
negative real parts, and $(a_1,a_2)$ generically unfold
the cusp.

\item Along this cusp continuation, for every sufficiently
small $\alpha,\beta,\gamma>0$, the system has
\emph{at least eight distinct locally exponentially stable
equilibria}: the origin, $E_{12}$, $E_{23}$, and at least
five strictly positive equilibria. These eight stable
equilibria are distinct from the interior cusp equilibrium.
Both proper siphon faces are transversely stable at
$E_{12}$ and $E_{23}$, respectively. No assertion of
exactly eight stable equilibria is made.
\end{enumerate}
\eeT

\begin{proof}
For the siphon classification, on $x_1=0$ the first
production argument equals $\alpha x_2$, so a siphon
containing $1$ must contain $2$. Similarly, one containing
$3$ must contain $2$. On $x_2=0$ the second argument equals
$\beta x_1x_3$, so a siphon containing $2$ must contain
at least one of $1$ and $3$. This gives exactly the three
listed siphons. At an equilibrium, $x_2>0$ implies
$x_1,x_3>0$ by the first and third equations; whereas
$x_2=0$ implies $x_1x_3=0$ by the second equation.
Thus the asserted supports are exhaustive. In the strictly
positive interior, the four cross derivatives associated
with the edges $2\to1$, $1\to2$, $3\to2$, $2\to3$ are
strictly positive. Hence the Jacobian is irreducible Metzler.
The box is invariant because $\Phi_a(q)\ge0$ for $q\ge0$ and
$\Phi_a(q)<1$ for every finite $q\ge0$.

On $\Sigma_{12}$ the sole resident coordinate satisfies
$\dot x_3=-x_3+\Phi_a(x_3)$, independently of the couplings;
the same scalar equation holds for $x_1$ on $\Sigma_{23}$.
Differentiating the absent equations gives
\eqref{eq:KBoverlapJtra}. Since the linear removal matrix
is the identity, the nonnegative production blocks have
Perron eigenvalues
$a_1(1+\sqrt{\alpha\beta s_a})$ and
$a_1(1+\sqrt{\gamma\beta s_a})$, respectively.
The tangential eigenvalue is $\lambda_s(a)<0$.
The triangular face Jacobians then give the stability
equivalences. At the origin, the $x_1x_3$ term has no
linear part, and the Jacobian is triangular after reordering;
all its eigenvalues equal $a_1-1<0$.

For cusp continuation, regard $(x_2,x_3)$ as the
hyperbolic variables and $x_1$ as the scalar critical
variable. At $(a,\alpha,\beta,\gamma)=(a^0,0,0,0)$
and $(x_1,x_2,x_3)=(r,s,s)$, the derivatives of the last
two equilibrium equations with respect to $(x_2,x_3)$
form $\lambda_s(a^0)I_2$, which is invertible.
The implicit-function theorem therefore solves those
equilibrium equations locally as
\[
 (x_2,x_3)=Y(x_1;a,\alpha,\beta,\gamma).
\]
Substitution in the first equilibrium equation produces
the scalar Lyapunov--Schmidt equation
\[
 g(x_1;a,\alpha,\beta,\gamma)
 =-x_1+\Phi_a\bigl(x_1+\alpha Y_1(x_1;a,\alpha,\beta,\gamma)\bigr).
\]
At zero coupling, $g(x_1)=\Phi_a(x_1)-x_1$. By
Proposition~\ref{p:Qua}, at $(x_1,a)=(r,a^0)$,
\[
 g=g_{x_1}=g_{x_1x_1}=0,\qquad
 g_{x_1x_1x_1}=\frac{36}{49}\ne0,
\]
and
\[
 \det\frac{\partial(g,g_{x_1})}{\partial(a_1,a_2)}
 =r^2(1-r)^4>0.
\]
At zero coupling $J=\operatorname{Diag}(0,\lambda_s(a^0),\lambda_s(a^0))$ and $f_1$ is independent of
$(x_2,x_3)$. Both null vectors are $e_1$ and $h_{20}=0$, so these are the cusp conditions of
Proposition~\ref{p:sadCus}(2). Proposition~\ref{p:sadCus}(3) gives the stated smooth cusp
continuation in $(\alpha,\beta,\gamma)$, with the cubic, unfolding, and stable
complementary-eigenvalue conditions. Strict positivity of the cusp coordinates and of
all four odds coefficients persists by continuity.

Finally we prove the stable-equilibrium count without
assuming that the nonlinear equations can be solved
explicitly. Let $S\subset\{1,2,3\}$ and write $z^S(a)$ for
the vector having coordinate $s_a$ at each index in $S$
and zero elsewhere. At zero coupling the eight $z^S(a)$
are distinct, hyperbolically stable equilibria of the
three decoupled scalar equations. Define the increasing
map $T_{a,\alpha,\beta,\gamma}:[0,1]^3\to[0,1]^3$ by
\[
 T(x)=\bigl(\Phi_a(x_1+\alpha x_2),
             \Phi_a(x_2+\beta x_1x_3),
             \Phi_a(x_3+\gamma x_2)\bigr).
\]
For every $S$, nonnegativity of the couplings gives
$T(z^S(a))\ge z^S(a)$ componentwise.
Since $\Phi_a'(0+)<1$ and $\Phi_a'(s_a)<1$, one can choose a
single small $h>0$, uniformly for $a$ near $a^0$, such
that the eight boxes
\[
 I_S(a)=[z^S(a),z^S(a)+h\mathbf1]
\]
are pairwise disjoint, lie in $[0,1]^3$, and satisfy
$\Phi_a(t+h)<t+h$ for $t\in\{0,s_a\}$.
For sufficiently small nonnegative couplings,
$T(z^S(a)+h\mathbf1)\le z^S(a)+h\mathbf1$.
Monotonicity then makes every $I_S(a)$ invariant under
$T$. By reducing $h$ and the coupling bound if necessary,
$\Phi_a'$ is uniformly less than one at all arguments attained
in these boxes, and the row sums of $D T$ are uniformly
less than one, since the row sums of the Jacobian of the arguments tend to one as
the couplings vanish. Therefore $T$ is a contraction on
each box. Each box contains a unique equilibrium
$e^S$, and $J(e^S)=DT(e^S)-I$ is Hurwitz.

The three boxes indexed by $S=\varnothing,\{1\},\{3\}$
contain the exact stable equilibria $E_0,E_{23},E_{12}$.
For the remaining five subsets, the equilibrium is
strictly positive: if $2\in S$ then $x_2>0$ forces
$x_1,x_3>0$; if $S=\{1,3\}$ then $x_1,x_3>0$
forces $x_2>0$. Thus at least five stable equilibria
are interior. Their boxes are separated from the cusp,
whose first coordinate is close to $r$, rather than
to $0$ or $s$. Along the cusp continuation all parameters
remain in the neighborhoods just used. This completes
the proof.
\end{proof}

\paragraph{Method: certify cusp continuation and simultaneous stability.}
Choose a square augmented system $\mathcal G(X;\eta)=0$ for an equilibrium, its simple zero
eigenvalue, the vanishing quadratic coefficient, and two adjustable unfolding parameters. Here $X$
collects the unknown state, normalized nullvectors and adjustable parameters; $\eta$ denotes the
remaining couplings. The standard cusp formulas already used in the paper supply these equations.
State the normalization explicitly and check that the augmented Jacobian is invertible at the seed.
For a box $\mathcal X$, centre $X_0$, coupling box $E$, and a fixed numerical inverse approximation
$C$, enclose
\[
 \mathcal K(\mathcal X,E)
 =X_0-C\mathcal G(X_0;E)
  +\bigl(I-C D_X\mathcal G(\mathcal X;E)\bigr)(\mathcal X-X_0).
\]
Certify strict inclusion in $\operatorname{int}\mathcal X$ and a uniform norm bound
\[
 \sup_{X\in\mathcal X,\,\eta\in E}\|I-C D_X\mathcal G(X;\eta)\|<1.
\]
With $C$ nonsingular, these conditions give a unique augmented solution in $\mathcal X$ for each
fixed $\eta\in E$. Separately certify positivity of the required coefficients and states,
nonvanishing of the cubic coefficient and unfolding determinant, and negativity of the real parts
of all complementary eigenvalues. Enclose the continued stable equilibria in disjoint boxes with
validated existence and Hurwitz tests, and enclose the resident equilibria and their transverse
invasion inequalities. A successful cusp solve alone certifies none of these additional assertions.

\beC[Three-population odds networks admit nine stable equilibria near the interior cusp][nine stable equilibria][Nin]
Fix any sufficiently small positive coupling triple
$(\alpha,\beta,\gamma)$ covered by Theorem~\ref{t:KBoverlapCusp}, and let
$a^c=(a_1^c,a_2^c)$ be its continued cusp coefficients, with
$a_3=128/343$ and $a_4=1024/343$ fixed.
There exists a nonempty open set of admissible coefficient pairs
$(a_1,a_2)$ whose closure contains $a^c$, on which the network has
at least nine distinct locally exponentially stable equilibria:
the origin, two equilibria on the proper siphon faces, and at least
six strictly positive equilibria.
Two additional strictly positive equilibria near the former cusp
each have exactly one unstable eigenvalue.
The minimal siphons remain $\Sigma_{12}=\{1,2\}$ and
$\Sigma_{23}=\{2,3\}$, and the two continued stable boundary equilibria
remain transversely stable.
\eeC
\begin{proof}
At $a=a^c$, Theorem~\ref{t:KBoverlapCusp} supplies eight distinct hyperbolically
stable equilibria, separated from the interior cusp equilibrium.
The implicit-function theorem continues these equilibria to a common
neighbourhood of $a^c$, preserving their distinctness and stability.
The origin remains an equilibrium. The two stable boundary equilibria
remain on their respective invariant siphon faces, and the other five
remain strictly positive.

Write the vector field as
$f=(f_1,f_2,f_3)^\top$ and put $H=(f_2,f_3)^\top$.
As in the proof of Theorem~\ref{t:KBoverlapCusp}, the last two equilibrium
equations can be solved near the cusp as
\[
 (x_2,x_3)=Y(x_1;a),
 \qquad
 g(x_1;a)=f_1(x_1,Y(x_1;a);a).
\]
At the cusp,
\[
 g=g_{x_1}=g_{x_1x_1}=0,\qquad
 g_{x_1x_1x_1}>0,\qquad
 \det\frac{\partial(g,g_{x_1})}{\partial(a_1,a_2)}\ne0.
\]
The positivity of the cubic follows by continuity from its reference
value $36/49$. The generic scalar cusp unfolding therefore has an open
parameter region, accumulating at $a^c$, with three simple nearby
equilibria. In their increasing $x_1$ order, the derivatives
$g_{x_1}$ are positive, negative, and positive, respectively.

To transfer these signs to the full system, let
\[
 J_{\mathrm{res}}
   =D_{(x_2,x_3)}H
\]
at a nearby equilibrium. This matrix remains Hurwitz, since it is close
to the strictly stable resident matrix used in the cusp construction.
The Schur-complement identity gives
\[
 \det D_xf=\det J_{\mathrm{res}}\,g_{x_1}.
\]
The two eigenvalues complementary to the critical eigenvalue remain
in the open left half-plane. Their product is positive, as is
$\det J_{\mathrm{res}}$. Hence the remaining real eigenvalue has the
same sign as $g_{x_1}$. The middle equilibrium is locally exponentially
stable, whereas each outer equilibrium has exactly one unstable
eigenvalue.

All three equilibria are strictly positive and lie in a neighbourhood
disjoint from the eight continued stable equilibria. The middle one
therefore raises the stable count to at least nine, including at least
six strictly positive equilibria. Shrinking the parameter region
preserves positivity of the odds coefficients and the strict
inequalities $R_{\Sigma_{12}}<1$ and $R_{\Sigma_{23}}<1$ at the two
continued boundary equilibria. The fixed positive couplings and
unchanged production supports preserve the minimal siphons.
\end{proof}

\beP[A composite energy certifies an explicit subset of the basin of a boundary equilibrium][][ComBas]
For \eqref{eq:KBoverlapCusp} with $\alpha,\beta,\gamma>0$ and $a_1>0$, let $\mathcal E_a$ be the
scalar energy following Proposition~\ref{p:Qua}. Take $\eta>0$, and put
$\ell=(\sqrt\beta,\sqrt\alpha)$, $\bar z=\eta/\min\{\sqrt\alpha,\sqrt\beta\}$,
$\bar q=\bar z(1+\max\{\alpha,\beta\})$ and $\rho=\sup_{0<q\le\bar q}\Phi_a(q)/q$. Assume
\[
 c:=1-\rho\bigl(1+\sqrt{\alpha\beta}\bigr)>0,
\]
and that $0<y_-<s_a<y_+<1$ satisfy $\Phi_a(q)>q$ on $[y_-,s_a)$, $\Phi_a(q)<q$ on $(s_a,y_+]$,
$\Phi_a(y_++\gamma\bar z)\le y_+$ and $\Phi_a'(s_a)<1$. Put
\[
 M_U=\max_{y\in[y_-,y_+]}|\Phi_a^{-1}(y)-y|,\qquad
 M_\Phi=\max_{q\in[y_-,\,y_++\gamma\bar z]}\Phi_a'(q).
\]
Then
\[
 K=\{x:\ x_1,x_2\ge0,\ \sqrt\beta\,x_1+\sqrt\alpha\,x_2\le\eta,\ y_-\le x_3\le y_+\}
\]
is forward invariant, and for $\kappa>\gamma M_UM_\Phi/(c\sqrt\alpha)$,
\[
 H(x)=\mathcal E_a(x_3)-\mathcal E_a(s_a)+\kappa\bigl(\sqrt\beta\,x_1+\sqrt\alpha\,x_2\bigr)
\]
is a strict Lyapunov function on $K$. Thus $E_{12}=(0,0,s_a)$ is globally asymptotically stable
relative to $K$, so $K$ is an explicit subset of its basin of attraction. Exchanging $x_1$ with $x_3$ and $\alpha$ with $\gamma$ gives the same statement for
$E_{23}$.
\eeP
\begin{proof}
Put $z=(x_1,x_2)$. On $K$, $x_1,x_2\le\bar z$ and $x_3\le1$, so the arguments of the first two equations lie
in $[0,\bar q]$. Hence $\dot x_1\le-x_1+\rho(x_1+\alpha x_2)$ and $\dot x_2\le-x_2+\rho(x_2+\beta x_1)$, and
therefore $\ell^\top\dot z\le-c\,\ell^\top z$. The faces $x_1=0$, $x_2=0$ point inward, and
$\ell^\top z$ decreases at the level $\eta$. On $x_3=y_-$, $\dot x_3\ge-y_-+\Phi_a(y_-)>0$; on $x_3=y_+$,
$\dot x_3\le-y_++\Phi_a(y_++\gamma\bar z)\le0$. So $K$ is forward invariant.

Let $U(y)=\mathcal E_a(y)-\mathcal E_a(s_a)$, so $U'(y)=\Phi_a^{-1}(y)-y$. The sign conditions and
monotonicity of $\Phi_a$ give $U'<0$ on $[y_-,s_a)$ and $U'>0$ on $(s_a,y_+]$; thus $U\ge0$ vanishes only at
$s_a$. Next, $U'(y)(-y+\Phi_a(y))=-(\Phi_a^{-1}(y)-y)(y-\Phi_a(y))$, and the quotient of
$(\Phi_a^{-1}(y)-y)(y-\Phi_a(y))$ by $(y-s_a)^2$ is continuous and positive on $[y_-,y_+]$, with limit
$(1-\Phi_a'(s_a))^2/\Phi_a'(s_a)$ at $s_a$. This gives the second inequality of
Proposition~\ref{p:ComLya}. Finally,
$U'(x_3)(\Phi_a(x_3+\gamma x_2)-\Phi_a(x_3))\le M_UM_\Phi\gamma x_2\le b\,\ell^\top z$, with
$b=\gamma M_UM_\Phi/\sqrt\alpha$. Proposition~\ref{p:ComLya} applies.
\end{proof}

The hypothesis $c>0$ is stronger than $R_{\Sigma_{12}}<1$, because it holds on all of $K$: it uses
$\rho\ge a_1$ and the bound $x_3\le1$ rather than the value $s_a$. At the couplings of
Table~\ref{tab:KBoverlapNum}, with the reference coefficients $(a_1,a_2)=(342/343,368/343)$,
$s_a=1/2$, $\bar z=1/10$, $y_-=1/5$ and $y_+=7/10$, one has $\rho\approx0.99998$ and
$c\approx1.6\cdot10^{-5}$; this numerical check, with the derivative bound tested at sampled points,
is illustrative.

\beP[A Lyapunov--Volterra function certifies attraction to a boundary resident]
[][LVboundary]
Under the hypotheses of Proposition~\ref{p:ComBas}, retain its invariant set $K$ and constants $c>0$, $M_\Phi$, $y_-$, $y_+$ and $s_a$. Define
\[
 U_V(y)=y-s_a-s_a\log\frac{y}{s_a},\qquad
 M_V=\max_{y\in[y_-,y_+]}\left|1-\frac{s_a}{y}\right|,\qquad
 b_V=\frac{\gamma M_\Phi M_V}{\sqrt{\alpha}}.
\]
For every $\kappa>b_V/c$, the function
\[
 H_V(x)=x_3-s_a-s_a\log\frac{x_3}{s_a}+\kappa\bigl(\sqrt{\beta}\,x_1+\sqrt{\alpha}\,x_2\bigr)
\]
is a strict Lyapunov function on $K$ for $E_{12}=(0,0,s_a)$. More precisely, put
\[
 a_*=\min_{y\in[y_-,y_+]}\frac{y-\Phi_a(y)}{y-s_a}>0,
\]
where the quotient at $y=s_a$ is defined as $1-\Phi_a'(s_a)$. Then
\[
 \dot H_V\le-\frac{a_*}{y_+}(x_3-s_a)^2-(\kappa c-b_V)\bigl(\sqrt{\beta}\,x_1+\sqrt{\alpha}\,x_2\bigr),
\]
and
\[
 H_V(x(t))\le e^{-\omega_Vt}H_V(x(0)),\qquad
 \omega_V=\min\left\{\frac{2a_*y_-^2}{s_ay_+},\ c-\frac{b_V}{\kappa}\right\}>0.
\]
Consequently every trajectory starting in $K$ converges to $E_{12}$.
\eeP
\begin{proof}
Write $y=x_3$ and $Z=\sqrt{\beta}\,x_1+\sqrt{\alpha}\,x_2$. The proof of Proposition~\ref{p:ComBas} gives $\dot Z\le-cZ$. For the resident equation,
\[
 U_V'(y)\bigl(-y+\Phi_a(y)\bigr)=-\frac{y-s_a}{y}\bigl(y-\Phi_a(y)\bigr)\le-\frac{a_*}{y_+}(y-s_a)^2.
\]
The additional production caused by $x_2$ satisfies
\[
 U_V'(y)\bigl[\Phi_a(y+\gamma x_2)-\Phi_a(y)\bigr]\le M_VM_\Phi\gamma x_2\le b_VZ.
\]
Adding $\kappa\dot Z$ proves the derivative inequality. The sign assumptions of Proposition~\ref{p:ComBas} and $\Phi_a'(s_a)<1$ make the quotient defining $a_*$ continuous and strictly positive.

Moreover, $U_V(s_a)=U_V'(s_a)=0$ and $U_V''(y)=s_a/y^2$, so
\[
 \frac{s_a}{2y_+^2}(y-s_a)^2\le U_V(y)\le\frac{s_a}{2y_-^2}(y-s_a)^2.
\]
These bounds give $\dot H_V\le-\omega_VH_V$ and the asserted convergence.
\end{proof}

This is the classical weighted Lyapunov--Volterra form about $E_{12}$. The positive resident coordinate contributes $x_3-s_a-s_a\log(x_3/s_a)$, and the coordinates whose equilibrium values are zero contribute linear terms, using $\lim_{p\downarrow0}(x-p-p\log(x/p))=x$. Thus the Fenchel--Bregman interpretation of Subsection~\ref{s:FenBre} has a direct application to the threshold network, without changing its equations or introducing an inverse response into the certificate. At the values above, $M_V=3/2$, $b_V\approx1.69\cdot10^{-4}$ and $b_V/c\approx10.4$ (illustrative).

For Theorem~\ref{t:Net} at $a=a^0$, suppose $0<\bar x<1/8$ and
$c_2:=1-(1+\gamma)\sup_{0<q\le\bar x(1+\gamma)}\Phi_{a^0}(q)/q>0$. Then for every $\varepsilon\ge0$ every
solution starting in $[0,\bar x]\times(0,1]$ converges to $(0,y^*)$. Indeed, $y\le1$ gives
$\dot x\le-c_2x$ on $[0,\bar x]$, so $x$ decays exponentially. Monotonicity of the response gives
$\dot y\ge-y+\Phi_{a^0}(10y)$, so comparison with the scalar resident equation, whose only positive
equilibrium is $y^*$, gives $\liminf y\ge y^*$. Since $y\le1$ and $x\to0$, the difference
$\Phi_{a^0}(y(10+\varepsilon x))-\Phi_{a^0}(10y)$ tends to zero uniformly, and the scalar field is negative
above $y^*$, so $\limsup y\le y^*$. On compact subsets of $[0,\bar x]\times(0,1)$ that are forward
invariant, the resident storage $\int\Psi^{-1}-y^2/2$ about $y^*$, with $\Psi(y)=\Phi_{a^0}(10y)$, plus a
linear term $\kappa x$ gives a Lyapunov function as in Proposition~\ref{p:ComLya}.

\beR[A double boundary transcritical relay]
[double boundary transcritical relay][KBrelay]
Fix $a$ with $a_1<1$ and $\alpha,\gamma>0$, and vary only $\beta$. Species $2$ belongs to both $\Sigma_{12}$ and
$\Sigma_{23}$, so its production, which contains $\beta x_1x_3$, gives a row of both transverse matrices. Since
$x_2=0$ at $E_{12}$ and $E_{23}$, the residents do not depend on $\beta$. Hence both invasion numbers increase
with $\beta$ at fixed residents, and
\[
 R_{\Sigma_{12}}=1\iff\beta=\beta_{12}:=\frac{(1-a_1)^2}{a_1^2\alpha s_a},\qquad
 R_{\Sigma_{23}}=1\iff\beta=\beta_{23}:=\frac{(1-a_1)^2}{a_1^2\gamma s_a}.
\]
Since the tangential eigenvalue $\lambda_s(a)$ is negative, each resident is stable below its value and unstable
above it. We call this a \emph{double boundary transcritical relay}, extending the boundary transcritical relays of
\cite{AAHrel}: one parameter takes two residents, on the faces of two overlapping siphons, through transcritical
bifurcations in turn.

Each crossing is a nondegenerate transcritical bifurcation when $a_2>a_1^2$. The zero eigenvalue is simple, its
irreducible Metzler block has Perron vectors $w,v\gg0$, and the parameter derivative
$v^\top(\partial_\beta J^\perp)w$ is positive. Since $\Phi_a(q)=a_1q+(a_2-a_1^2)q^2+O(q^3)$, all second-order
terms along the critical direction are nonnegative, with a positive sum, so the positive branch
exists for $\beta$ below the critical value: the bifurcation is backward. Along the branch from $E_{12}$, the second equation gives
\[
 \beta=\frac{\Phi_a^{-1}(x_2)-x_2}{x_1x_3},
\]
with $x_1,x_3$ the roots of the first and third equations near $0$ and $s_a$. Figure~\ref{fig:relay} is computed
from this formula.

\emph{Numerical evidence.} The following was computed for the reference coefficients $a=(342,368,128,1024)/343$, so
$s_a=1/2$ and $a_2-a_1^2\approx0.079$, with $\alpha=2\cdot10^{-7}$, $\gamma=5\cdot10^{-8}$, which give $\beta_{12}\approx85.5$ and
$\beta_{23}\approx342$. Trajectories were integrated to $t=5\cdot10^7$ from $x_2(0)=10^{-4}$.
\begin{enumerate}
\item At $E_{12}$, with the critical Perron vectors normalized to have second coordinate one,
$v^\top(\partial_\beta J)w\approx3.41\cdot10^{-5}$ and $\frac12v^\top D^2f[w,w]\approx0.079$. The predicted branch agrees with Newton
solves: at $\beta=0.999\,\beta_{12}$ it predicts $x_2\approx3.682\cdot10^{-5}$, and the solve gives $3.684\cdot10^{-5}$.
\item The largest real part of the Jacobian spectrum lies in $[3.3\cdot10^{-6},1.04\cdot10^{-3}]$ along the sampled backward branches,
which are therefore unstable, and in $[-2.9\cdot10^{-3},-2.6\cdot10^{-3}]$ along the upper branches, which are stable.
\item Below its value each resident attracts the perturbed state: the trajectories return to $E_{12}$ at $0.8\,\beta_{12}$ and to $E_{23}$
at $0.8\,\beta_{23}$. Past its value the state jumps to an upper branch: the end states are $(3.66\cdot10^{-5},0.5343,0.5000)$ at
$1.2\,\beta_{12}$ from $E_{12}$, and $(0.5000,0.5343,9.1\cdot10^{-6})$ at $1.2\,\beta_{23}$ from $E_{23}$.
\item Between the two values one resident has been lost while the other survives: at $\beta=1.2\,\beta_{12}$ the trajectory from $E_{23}$ returns to
$E_{23}$.
\item The jump is not undone by decreasing $\beta$, down to the smallest value tested: from the upper state reached at
$1.2\,\beta_{12}$, the trajectory stays on the upper branch at $0.5\,\beta_{12}$, with $x_2\approx0.516$, and at $0.05\,\beta_{12}$, with $x_2\approx0.502$.
This is consistent with the backward branch reaching $\beta=0$ near $x_2=1/8$, so that no fold closes a hysteresis loop for $\beta>0$.
\end{enumerate}

\emph{Status.} Established here are the formulas for $\beta_{12},\beta_{23}$, the stability change of the residents, and,
subject to checking the sketch above, nondegeneracy and the backward direction under $a_2>a_1^2$. Items (2)--(5) are computations at one choice of
coefficients and couplings, and are not proved in general. It is open whether the upper branch, its stability and the irreversibility hold
throughout the family, and whether \emph{relay} as defined in \cite{AAHrel}, where a crossing produces a positive branch on an adjacent
face, covers this case, in which the branch leaving the face is interior. Each single crossing is a standard boundary transcritical (invasion)
bifurcation. A limited literature search found no name for the feature emphasized here, that one production term of a shared
species enters two transverse matrices while both residents stay fixed, so the term is proposed provisionally and the feature is not claimed to be new.

The relay requires non-disjoint siphons. If $\Sigma$ and $\Sigma'$ are disjoint, the production of a species of $\Sigma$
gives no row of $F_{\Sigma'}$, so changing it can affect $J_{\Sigma'}^\perp$ only by moving the resident.
Folds and cusps do not distinguish the two cases: Theorem~\ref{t:Net} has them with the disjoint singleton siphons
$\{x\}$ and $\{y\}$. The values above concern the boundary residents, and do not extend the small-coupling
range of Theorem~\ref{t:KBoverlapCusp}: they lie far outside it, so the relay is not shown to occur next to the interior cusp.
\eeR

\begin{figure}[H]
\centering
\begin{minipage}{.48\linewidth}
\centering
\begin{tikzpicture}
\begin{axis}[width=\linewidth,height=.8\linewidth,xlabel={$\beta$},ylabel={$x_2$},xmin=0,xmax=450,ymin=-.03,ymax=.75,
 xtick={0,200,400},tick label style={font=\scriptsize},label style={font=\small},title style={font=\small},extra x ticks={85.5},extra x tick labels={$\beta_{12}$},title={from $E_{12}$}]
\addplot[blue,thick] coordinates {(0,0) (85.50,0)};
\addplot[blue,dashed] coordinates {(85.50,0) (450,0)};
\addplot[blue,dashed] coordinates {(85.50,0.0000) (85.03,0.0002) (84.79,0.0003) (84.43,0.0005) (83.88,0.0007) (83.04,0.0011) (81.79,0.0016) (79.91,0.0025) (77.12,0.0037) (73.01,0.0057) (67.07,0.0087) (58.74,0.0132) (47.56,0.0200) (40.44,0.0250) (29.80,0.0338) (21.34,0.0426) (14.76,0.0515) (9.77,0.0603) (6.12,0.0691) (3.55,0.0779) (1.85,0.0867) (0.82,0.0955) (0.27,0.1044) (0.05,0.1132) (0.00,0.1220) (0.00,0.1250)};
\addplot[blue,thick] coordinates {(0.00,0.5000) (4.85,0.5020) (24.61,0.5096) (46.37,0.5173) (70.22,0.5249) (96.25,0.5326) (124.58,0.5402) (155.28,0.5478) (188.48,0.5555) (224.28,0.5631) (262.80,0.5708) (304.15,0.5784) (348.45,0.5861) (395.85,0.5937) (446.47,0.6013) (455.00,0.6026)};
\draw[red,thick,-{Stealth[length=5pt]}] (axis cs:85.50,0) -- (axis cs:85.50,0.5295);
\end{axis}
\end{tikzpicture}
\end{minipage}\hfill
\begin{minipage}{.48\linewidth}
\centering
\begin{tikzpicture}
\begin{axis}[width=\linewidth,height=.8\linewidth,xlabel={$\beta$},ylabel={$x_2$},xmin=0,xmax=450,ymin=-.03,ymax=.75,
 xtick={0,200,400},tick label style={font=\scriptsize},label style={font=\small},title style={font=\small},extra x ticks={342.0},extra x tick labels={$\beta_{23}$},title={from $E_{23}$}]
\addplot[blue,thick] coordinates {(0,0) (341.99,0)};
\addplot[blue,dashed] coordinates {(341.99,0) (450,0)};
\addplot[blue,dashed] coordinates {(341.99,0.0000) (340.13,0.0002) (339.17,0.0003) (337.71,0.0005) (335.51,0.0007) (332.18,0.0011) (327.16,0.0016) (319.65,0.0025) (308.47,0.0037) (292.03,0.0057) (268.30,0.0087) (234.97,0.0132) (190.26,0.0200) (161.75,0.0250) (119.20,0.0338) (85.38,0.0426) (59.06,0.0515) (39.10,0.0603) (24.46,0.0691) (14.18,0.0779) (7.39,0.0867) (3.26,0.0955) (1.09,0.1044) (0.20,0.1132) (0.00,0.1220) (0.00,0.1250)};
\addplot[blue,thick] coordinates {(0.00,0.5000) (19.42,0.5020) (98.51,0.5096) (185.60,0.5173) (281.08,0.5249) (385.30,0.5326) (455.00,0.5373)};
\draw[red,thick,-{Stealth[length=5pt]}] (axis cs:341.99,0) -- (axis cs:341.99,0.5295);
\end{axis}
\end{tikzpicture}
\end{minipage}
\caption{The double boundary transcritical relay of Remark~\ref{r:KBrelay} for $\dot x_1=-x_1+\Phi_a(x_1+\alpha x_2)$, $\dot x_2=-x_2+\Phi_a(x_2+\beta x_1x_3)$, $\dot x_3=-x_3+\Phi_a(x_3+\gamma x_2)$: equilibria, with the backward branches given by $\beta=(\Phi_a^{-1}(x_2)-x_2)/(x_1x_3)$, for the reference coefficients and the couplings
$\alpha,\gamma$ of Table~\ref{tab:KBoverlapNum}: the shared species $x_2$ against $\beta$. Solid curves are stable and dashed curves
unstable. Each resident ($x_2=0$) loses stability at its own value, $\beta_{12}\approx85.5$ or $\beta_{23}\approx342$; the unstable
positive branch leaves it backward and ends near $x_2=1/8$ as $\beta\to0$. Past its value the state jumps to the upper
branch, where it remains when $\beta$ is decreased.}
\label{fig:relay}
\end{figure}

\beR[The siphon lattice changes at zero coupling]
The proof starts from the decoupled system $\alpha=\beta=\gamma=0$, at which
every singleton $\{i\}$ is a siphon and every coordinate face is invariant. For
$\alpha,\beta,\gamma>0$ the minimal siphons are instead $\Sigma_{12}$ and
$\Sigma_{23}$. The implicit-function arguments only need the decoupled point as
a base for continuation, so this causes no gap; but the siphon lattice is not
constant across that point. Theorem~\ref{t:KBoverlapCusp} is a statement about
the positive-coupling family, which retains the overlapping-siphon lattice of
Example~\ref{e:KBoverlap}; it is not a continuation through zero coupling with an
unchanged lattice. Likewise, $R_{\Sigma_{12}}$ and $R_{\Sigma_{23}}$ depend on the
higher coefficients $a_2,a_3,a_4$ through $s_a$, although these coefficients do
not appear explicitly in \eqref{eq:oddsTransverse}.
\eeR

\paragraph{A numerical instance.}
Table~\ref{tab:KBoverlapNum} reports one coupling triple,
$(\alpha,\beta,\gamma)=(2\cdot10^{-7},10^{-6},5\cdot10^{-8})$. The cusp was
located by Newton's method on the seven equations $f=0$, $Jv=0$ ($v_1=1$),
$w^\top D^2f(v,v)=0$ with $w$ a left null vector, in 60-digit arithmetic.
The cubic coefficient and the unfolding determinant are those of the
Lyapunov--Schmidt function $g$ of the proof. They were computed by central
differences and agree, to the digits shown, with implicit differentiation of
the resident equations. The eight equilibria were found by Newton's method from
the points $z^S(a)$. The computation is illustrative, not interval-validated;
the Mathematica script \texttt{overlNum.wl}, distributed with the source,
reproduces the full-precision values.
\begin{table}[H]
\centering\small
\begin{tabular}{ll}
\hline
Quantity & Value\\\hline
continued coefficients $(a_1,a_2)$ & $(0.997082757,\ 1.072892269)$\\
cusp $(x_1,x_2,x_3)$ & $(0.124986385,\ 0.500005319,\ 0.500004404)$\\
complementary eigenvalues & $-0.0393616,\ -0.0393609$\\
$g_{x_1x_1x_1}$ (reference $36/49\approx0.734694$) & $0.734463$\\
$\det\partial(g,g_{x_1})/\partial(a_1,a_2)$ (reference $0.00915909$) & $0.00915761$\\
resident level $s_a$ & $0.500003793$\\
$R_{\Sigma_{12}}$, $R_{\Sigma_{23}}$ & $0.997083072,\ 0.997082914$\\\hline
stable equilibrium, box $S$ & spectral abscissa\\\hline
$\varnothing$: $(0,0,0)$ & $-0.002917$\\
$\{3\}$: $(0,0,0.500003793)=E_{12}$ & $-0.002917$\\
$\{1\}$: $(0.500003793,0,0)=E_{23}$ & $-0.002917$\\
$\{2\}$: $(3.42\cdot10^{-5},\ 0.500003793,\ 8.55\cdot10^{-6})$ & $-0.002912$\\
$\{2,3\}$: $(3.42\cdot10^{-5},\ 0.500003794,\ 0.500004404)$ & $-0.002912$\\
$\{1,3\}$: $(0.500003794,\ 8.56\cdot10^{-5},\ 0.500003793)$ & $-0.002904$\\
$\{1,2\}$: $(0.500006234,\ 0.500003793,\ 8.55\cdot10^{-6})$ & $-0.002916$\\
$\{1,2,3\}$: $(0.500006234,\ 0.500009895,\ 0.500004404)$ & $-0.039361$\\\hline
\end{tabular}
\caption{Theorem~\ref{t:KBoverlapCusp} at
$(\alpha,\beta,\gamma)=(2\cdot10^{-7},10^{-6},5\cdot10^{-8})$, with $a_3,a_4$
fixed. Both invasion numbers are below one and differ, since $\alpha\ne\gamma$.}
\label{tab:KBoverlapNum}
\end{table}

\paragraph{A nearby instance with nine stable equilibria.}
Corollary~\ref{c:Nin} concerns coefficient pairs near, but not at, the cusp.
Keeping the couplings of Table~\ref{tab:KBoverlapNum}, the same script moves
$(a_1,a_2)$ from the cusp values by the solution $\delta a$ of
$\partial_a(g,g_{x_1})\,\delta a=(0,-10^{-6})$, which makes $g_{x_1}$ negative at
the cusp point. Newton's method from the eight points $z^S(b)$ and from three
points near the former cusp gives eleven distinct equilibria, with residuals
below $10^{-52}$ and minimal pairwise distance $2.85\cdot10^{-3}$
(Table~\ref{tab:KBnineNum}). Nine of them have negative spectral abscissa. The
middle equilibrium near the former cusp is stable, with spectral abscissa
$-1.0\cdot10^{-6}$, of the order of the shift in $g_{x_1}$; the two outer ones
have one positive eigenvalue each. This is again illustrative, not
interval-validated.
\begin{table}[H]
\centering\small
\begin{tabular}{ll}
\hline
Quantity & Value\\\hline
coefficients $(a_1,a_2)$ & $(0.997084063,\ 1.072881819)$\\
resident level $s_b$ & $0.499991347$\\
$R_{\Sigma_{12}}$, $R_{\Sigma_{23}}$ & $0.997084378,\ 0.997084221$\\\hline
equilibrium & spectral abscissa\\\hline
the eight continued equilibria of Table~\ref{tab:KBoverlapNum} & between $-0.03935$ and $-0.002902$\\
$(0.122134360,\ 0.499992838,\ 0.499991957)$ & $+1.97\cdot10^{-6}$\\
$(0.124986385,\ 0.499992872,\ 0.499991957)$ & $-1.00\cdot10^{-6}$\\
$(0.127850842,\ 0.499992907,\ 0.499991957)$ & $+2.03\cdot10^{-6}$\\\hline
\end{tabular}
\caption{Corollary~\ref{c:Nin} at the couplings of Table~\ref{tab:KBoverlapNum}:
nine locally exponentially stable equilibria, the ninth being the middle
equilibrium of the cusp unfolding.}
\label{tab:KBnineNum}
\end{table}

\paragraph{Interpretation.}
The bare model of Example~\ref{e:KBoverlap} has overlapping minimal
siphons but no positive boundary resident; its self-feedback
augmentation \eqref{eq:KBoverlapCusp} retains the same siphon lattice.
Weak-coupling persistence of combinatorial families of stable
equilibria is already in \cite{Boud26}, so the count of eight is
not by itself the point; nor is the ninth stable equilibrium of
Corollary~\ref{c:Nin}, which exists only off the cusp parameter and relies on
the positive cubic coefficient. What the theorem adds is the combination:
the stable equilibria coexist with a nondegenerate interior cusp,
the coupling preserves two overlapping minimal siphons, the proper
siphon faces carry nonzero resident equilibria, and distinct,
explicit invasion numbers decide their transverse stability.
This is local coexistence; the theorem does not assert that the cusp
itself crosses either invasion-threshold hypersurface as a coupling is varied.